\section{Ringed spaces}
\label{section-ringed-spaces}

\noindent
Let $X$ be a topological space and let $\mathcal{O}_X$
be a sheaf of rings on $X$. We are supposed to
think of the sheaf of rings $\mathcal{O}_X$ as a sheaf of functions
on $X$. And if $f : X \to Y$ is a ``suitable'' map, then by composition
a function on $Y$ turns into a function on $X$.
Thus there should be a natural
$f$-map from $\mathcal{O}_Y$ to $\mathcal{O}_X$,
see Definition \ref{definition-f-map} and
Lemma \ref{lemma-f-map}.
For a precise example, see
Example \ref{example-continuous-map-ringed} below. Here is the
relevant abstract definition.

\begin{definition}
\label{definition-ringed-space}
A {\it ringed space} is a pair $(X, \mathcal{O}_X)$ consisting
of a topological space $X$ and a sheaf of rings $\mathcal{O}_X$
on $X$. A {\it morphism of ringed spaces}
$(X, \mathcal{O}_X) \to (Y, \mathcal{O}_Y)$ is a pair
consisting of a continuous map $f : X \to Y$ and an
$f$-map of sheaves of rings
$f^\sharp : \mathcal{O}_Y \to \mathcal{O}_X$.
\end{definition}

\begin{example}
\label{example-continuous-map-ringed}
Let $f : X \to Y$ be a continuous map of topological spaces.
Consider the sheaves of continuous real valued functions
$\mathcal{C}^0_X$ on $X$ and $\mathcal{C}^0_Y$ on $Y$,
see Example \ref{example-C0-sheaf-rings}. We claim that there
is a natural $f$-map $f^\sharp : \mathcal{C}^0_Y \to \mathcal{C}^0_X$
associated to $f$. Namely, we simply define it by the
rule
\begin{eqnarray*}
\mathcal{C}^0_Y(V) & \longrightarrow & \mathcal{C}^0_X(f^{-1}V) \\
h & \longmapsto & h \circ f
\end{eqnarray*}
Strictly speaking we should write $f^\sharp(h) = h \circ f|_{f^{-1}(V)}$.
It is clear that this is a family of maps as in
Definition \ref{definition-f-map} and compatible with the $\mathbf{R}$-algebra
structures. Hence it is an $f$-map of sheaves of $\mathbf{R}$-algebras, see
Lemma \ref{lemma-f-map-sets-algebraic-structures}.

\medskip\noindent
Of course there are lots of other situations where there is
a canonical morphism of ringed spaces associated to a geometrical
type of morphism. For example, if $M$, $N$ are $\mathcal{C}^\infty$-manifolds
and $f : M \to N$ is an infinitely differentiable map, then
$f$ induces a canonical morphism of ringed spaces
$(M, \mathcal{C}_M^\infty) \to (N, \mathcal{C}^\infty_N)$.
The construction (which is identical to the above) is left
to the reader.
\end{example}

\noindent
It may not be completely obvious how to compose morphisms of
ringed spaces hence we spell it out here.

\begin{definition}
\label{definition-composition-maps-ringed-spaces}
Let
$(f, f^\sharp) : (X, \mathcal{O}_X) \to (Y, \mathcal{O}_Y)$ and
$(g, g^\sharp) : (Y, \mathcal{O}_Y) \to (Z, \mathcal{O}_Z)$
be morphisms of ringed spaces. Then we define
the {\it composition of morphisms of ringed spaces}
by the rule
$$
(g, g^\sharp) \circ (f, f^\sharp) = (g \circ f, f^\sharp \circ g^\sharp).
$$
Here we use composition of $f$-maps defined in
Definition \ref{definition-composition-f-maps}.
\end{definition}


